\def\ChapterCode{RSP}
\chapter
[\CO{[ATM]} Störungen der Atemwege und der Atmung]
{Störungen der Atemwege und der Atmung}
\label{topic:respiratorische-notfaelle}
\label{chp:atm}
\ChapterIntro
{%
~~~
}
{%
\begin{description}
  \item[Maintainer:] Sebastian Gabriel
  \item[Autoren:] Diverse
  \item[Reviewer:] Standard-Reviewprozess
  \item[Version:] {\VarDocVersionTypeName} ({\VarDocVersionTypeDescription})
  \item[SHA1:] {\DocGitCommit}
\end{description}

{\TxTeilkapitelAASS}
}

% % 1-E
% \TextSlide
% {\TxABCDE}
% {~
% \begin{ABCDEText}
%   \item[\ABCDE{1}] 
%   \item[\ABCDE{2}] 
%   \item[\ABCDE{3}] 
%   \item[\ABCDE{4}] 
%   \item[\ABCDE{A}] 
%   \item[\ABCDE{B}] 
%   \item[\ABCDE{C}] 
%   \item[\ABCDE{D}] 
%   \item[\ABCDE{E}] 
%   \item[\ABCDE{=}] 
%   \item[\ABCDE{\ldots}] 
% \end{ABCDEText}
% }
% {
% \begin{ABCDESlide}
%   \item[\ABCDE{Sz}] 
%   \item[\ABCDE{Ei}] 
%   \item[\ABCDE{Be}] 
%   \item[\ABCDE{Hb}] 
%   \item[\ABCDE{A}] 
%   \item[\ABCDE{B}] 
%   \item[\ABCDE{C}] 
%   \item[\ABCDE{D}] 
%   \item[\ABCDE{E}] 
%   \item[\ABCDE{=}] 
%   \item[\ABCDE{\ldots}] 
% \end{ABCDESlide}
% }
% {}
% {}
% {}
% 
% \TextSlide
% {{\TxSAMPLER}}
% {
% \begin{SAMPLERText}
%   \item[\SAMPLER{S}] 
%   \item[\SAMPLER{A}] 
%   \item[\SAMPLER{M}] 
%   \item[\SAMPLER{P}] 
%   \item[\SAMPLER{L}] 
%   \item[\SAMPLER{E}] 
%   \item[\SAMPLER{R}] 
% \end{SAMPLERText}
% }
% {
% \begin{SAMPLERSlide}
%   \item[\SAMPLER{S}] 
%   \item[\SAMPLER{A}] 
%   \item[\SAMPLER{M}] 
%   \item[\SAMPLER{P}] 
%   \item[\SAMPLER{L}] 
%   \item[\SAMPLER{E}] 
%   \item[\SAMPLER{R}] 
% \end{SAMPLERSlide}
% }
% {}
% {}
% {}

\clearpage
\Text{{\TxQuerverweise}}{%
\begin{itemize}
  \item Lungenentzündung (Pneumonie): \dotfill \FullPrettyRef{topic:pneumonie}
  \item Tuberkulose:
  \dotfill \FullPrettyRef{topic:tuberkulose}
\end{itemize}
}{}{}{}{}





% 
% \section{Blutungen au dem Nasenraum}
% 
% % TODO: Nasenbluten
% 
% \Text
% {{\TxBeschreibung}}
% {}
% {}
% {}
% {}
% {}
% 

%%%%%%%%%%%%%%%%%%%%%%%%%%%%%%%%%%%%%%%%%%%%%%%%%%%%%%%%%%%
%%%%%%%%%%%%%%%%%%%%%%%%%%%%%%%%%%%%%%%%%%%%%%%%%%%%%%%%%%%
%%%%%%%%%%%%%%%%%%%%%%%%%%%%%%%%%%%%%%%%%%%%%%%%%%%%%%%%%%%
\section{Allgemein: \HeadingKeyword{Störungen} der Atemwege 
und der Atmung}
\label{topic:atemstörungen}
\label{topic:atemstillstand}

\TextSlide{Einleitung}
{%
\index{Ateminsuffizienz}%
\index{Insuffizienz!Atem-}%
\index{Dyspnoe}%
\index{Apnoe}%
\index{Bradypnoe}%
\index{Tachypnoe}%
Störungen der Atmung können aus vielen Ursachen eintreten.
\FullPrettyRef{tab:atemstoerungen-ursachen}, beschreibt die
wichtigsten Pathomechanismen und Störungen.
Der Begriff \T{Ateminsuffizienz} beschreibt allgemein eine
\EE{nicht ausreichende Atmung}, unabhängig von der
Ursache der Störung. Dies führt zu einem \ce{O2}-Mangel
(Hypoxie). Oft hat der Patient das Gefühl nicht genug Luft
zu bekommen und leidet an \EE{Atemnot} (\T{Dyspnoe}).
Der \T{Atemstillstand} (\T{Apnoe}) ist die extremste Form
der Ateminsuffizenz. Störungen der Atemgeschwindigkeit
werden als \I{Bradypnoe} (zu langsam) und als \I{Tachypnoe}
(zu schnell) bezeichnet.
}{
\begin{itemize}
  \item Ateminsuffizienz \ldots{} nicht ausreichende Atmung
  \item Hypoxie \ldots{} \ce{O2}-Mangel
  \item Atemnot (Dyspnoe)
  \item Atemstillstand (Apnoe)
  \item Bradypnoe (zu langsam)
  \item Tachypnoe (zu schnell)
\end{itemize}
}{}{}{}


\TextSlide{Gefahr der Ateminsuffizienz}
{%
Bei einer Ateminsuffizienz sinkt der \EE{\ce{O2}-Gehalt} im
Blut. Es kommt zu einer \EE{Hypoxie} und in weiterer Folge
zu einer Unterversorgung lebenswichtiger Organe wie Herz und
Hirn.
Doch auch die \EE{\ce{CO2}-Abgabe} ist gestört und es kommt
zu einer Anhäufung von \ce{CO2} im Blut. In weiterer Folge
kann es zu einer \E{atmungsbedingten Übersäuerung}
(\EE{Respiratorische Azidose},
\FullPrettyRef{topic:azidose-respiratorische}) kommen.
}{
\begin{itemize}
    \item \ce{O2}-Gehalt im Blut ↓ (Hypoxie)
    \item \ce{CO2}-Gehalt im Blut ↑
    \item Respiratorische Azidose
    (\FullPrettyRef{topic:azidose-respiratorische})
\end{itemize}
}{}{}{}

\AuthNotes{GABG: pathologische Atmungstypen inkl.
Zeichnung einbinden? Kussmaul, etc.}



\TextSlide{Ursachen}
{%
\prettyref{tab:atemstoerungen-ursachen}.
}{\prettyref{tab:atemstoerungen-ursachen}.}{}{}{}

\Layout{57}
{}
{}
{}
{\begin{tabulary}{\linewidth}[H]{LL}\addlinespace\toprule\addlinespace
\noindent\TabHeading{Störung des/der}  &
\noindent\TabHeading{Vorkommen und Erkrankungen}
\\\addlinespace\midrule\addlinespace
\noindent\TabSub{\ce{O2}-Angebotes}
 &
Gärkeller (\ce{CO2}\,${\uparrow}$, \ce{O2}\,${\downarrow}$)

\enquote{dünne Luft} (z.\,B. Hochgebirge;
niedriger Sauerstoffpartialdruck)

Defekte Gastherme (\ce{CO}\,${\uparrow}$,
\ce{O2}\,${\downarrow}$)

Ertrinken
\\\addlinespace
\noindent\TabSub{\ce{O2}-Diffusion}
 &
Lungenödem

Lungenentzündung

Lungenembolie
\\\addlinespace
\noindent\TabSub{Atemmechanik}
 &
Verlegung der oberen Atemwege (Zunge, Erbrochenes, Laryngospasmus,
         Glottisödem, Bolus)

Verlegung der unteren Atemwege (Bronchitis, Asthma,
Lungenödem)

Verminderung der Dehnbarkeit der Thoraxwand oder des Lungengewebes
     (Rippenfraktur, Pneumothorax, Pleuraerguß, Emphysem)
\\\addlinespace
\noindent\TabSub{Neuro\-muskulären Atemregulation}
 &
SHT (Schädel Hirn Trauma)

Schlaganfall (Insult)

Vergiftungen

Tumore

Hohe Querschnittslähmung (Rückenmark durchtrennt)

Muskelerkrankungen

Nervenerkrankung

Stoffwechselstörungen
\\\addlinespace
\noindent\TabSub{Sonstige Regulationsstörung}
 &
 Psychogen (z.\,B. Panikattacke)
\\\addlinespace\bottomrule
\end{tabulary}}
{Ursachen für
Atemstörungen.\label{tab:atemstoerungen-ursachen}}
{}


\TextSlide
{{\eye} Befunde}
{Atemstörungen können viele verschiedene Symptome
verursachen. \FullPrettyRef{tab:atmung-symptome}, listet
verschiedene Symptome auf. Neben technischen Geräten ist vor
allem \E{sehen},
\E{hören}, \E{fühlen} und ein \E{aufmerksamer Blick}
gefordert.
Die Abgrenzung der Atemstörung zur Kreislaufstörung ist oft
nicht einfach. Eine Störung des Kreislaufes hat meist auch
einen verminderten Sauerstoff\E{transport} zur Folge, daher
kommt es dabei häufig {auch} zur Atemnot.

Häufig erkennbar sind z.\,B. eine \EE{Zyanose}, 
Klage über subjektive \EE{Atemnot}, 
eine zu schnelle
oder langsame Atmung (Tachy-/Bradypnoe; 
\EE{AF <\,8\PerMin* oder >\,30\PerMin*}), 
eine auffällige \EE{Atemtiefe} (flach, tief, evtl. 
\EE{Schnappatmung}), 
pathologische Atemmuster, 
eine paradoxe oder überhaupt fehlende Atmung oder 
andere Symptome der Atemnot
(\prettyref{tab:atmung-symptome}).
\Z{Bei der Auskultation kann eine ungleichseitige Belüftung der 
Lungen festgestellt werden, 
sowie diverse Atemnebengeräusche 
(feuchte Rasselgeräusche, spastische Atemgeräusche).}

Die \EE{Sauerstoffsättigung} kann vermindert 
oder hoch sein (Kohlenmonoxid, Hyperventilationssyndrom). 
Ebenso wichtig sind Umstände wie das Vorhandensein 
von Heimsauerstoff
oder anderen Atemhilfen.
}
{\begin{itemize}
  \item Zyanose, Atemnot
  \item Tachy-/Bradypnoe (AF <\,8\PerMin* oder >\,30\PerMin*)
  \item Atemtiefe (flach, tief, Schnappatmung)
  \item Path. Atemmuster
  \item Paradoxe, fehlende Atmung
  \item Sauerstoffsättigung
  \item \Z{Auskultation: Belüftung, feuchte
  Rasselgaräusche, Spastik}
  \item Heimsauerstoff, Atemhilfen
  \item \FullPrettyRef{tab:atmung-symptome}
\end{itemize}
}
{}
{}
{}


\Layout{57}
{}
{}
{}
{\begin{tabulary}{\linewidth}[H]{LLLL}
\toprule\addlinespace
\noindent\TabHeading{Kriterium}
 &
\noindent\TabHeading{Befund}
 &
\noindent\TabHeading{Beschreibung}
 &
%\noindent\TabHeading{Anm.}
\\\addlinespace\midrule\addlinespace
\noindent\TabSub{Beschwerde}
 &
Atemnot (Dyspnoe) \index{Dyspnoe}
 &
Leitsymptom
 &
\\\addlinespace
\noindent\TabSub{Atemgeräusch}
 &
Brodeln \index{Atemgeräusch!brodelndes}
 &
Blubbern, klassisch für Lungenödem
\index{Atemgeräusch!blubberndes} &
\\\addlinespace
 &
Stridor \index{Stridor}
 &
 &
\\\addlinespace
 &
Brummen, Giemen
\index{Atemgeräusch!brummendes}\index{Atemgeräusch!giemendes}
& &
\\\addlinespace
 &
Rasselgeräusche \index{Rasselgeräusche}
 &
 &
\\\addlinespace
\noindent\TabSub{Frequenz}
 &
Beschleunigt
 &
Tachypnoe \index{Tachypnoe}
 &
\\\addlinespace
 &
Verlangsamt
 &
Bradypnoe \index{Bradypnoe}
 &
\\\addlinespace
\noindent\TabSub{Atemzugs\-volumen}
 &
Schnappatmung
 &
\FullPrettyRef{topic:schnappatmung}
 &
\\\addlinespace
 &
Flache Atmung
 &
 &
\\\addlinespace
 &
Tiefe Atmung
 &
Z.\,B. Kußmaul'sche Atmung
 &
\\\addlinespace
\noindent\TabSub{Hautfarbe}
 &
Blass
 &
Normal, Anämie, Blutverlust
 &
\\\addlinespace
 &
Rosig
 &
Normal, \ce{CO}-Vergiftung
 &
\\\addlinespace
 &
Bläulich 
 &
Zyanose: Hypoxie!
 &
\\\addlinespace
\noindent\TabSub{Körperlich}
 &
Einziehungen an den Rippen
&
Einsatz der Atemhilfsmuskulatur
&
\\\addlinespace
&
Aufstützen
&
Einsatz der Atemhilfsmuskulatur
&
\\\addlinespace
&
Nasenflügeln \index{Nasenflügeln}
&
Atemnot, besonders bei Kleinkindern
&
\\\addlinespace
&
Aufrechte Position
&
&
\\\addlinespace
&
\T[Atmung!paradoxe]{Paradoxe Atmung}
&
Brustkorb \E{senkt} sich bei Einatmung, z.\,B. bei
Serienrippenfraktur 
&
\\\addlinespace
\noindent\TabSub{Geräte}
&
Pulsoxymetrie
&
\FullPrettyRef{topic:pulsoxymetrie}
&
\\\addlinespace
&
\Z{Kapnometrie}
&
\Z{\ce{CO2}-Messung}
\\\addlinespace\bottomrule
\end{tabulary}}
{Symptome: Atmung.\label{tab:atmung-symptome}}
{}





% \clearpage
\section{Mechanische Atemwegsverlegung}
\label{topic:mechanische-atemwegsverlegung}
\label{topic:fremdkoerperaspiration}

\Layout{60}{\TxBeschreibung}
{%
\index{Aspiration}%
Wenn ein Fremdkörper den Atemweg teilweise oder vollständig 
\E{blockiert} oder \E{behindert}, spricht man von einer
mechanische Verlegung. Abhängig von der Größe des
Fremdkörpers können sowohl die oberen als auch die unteren
Atemwege betroffen sein. Wenn der Fremdkörper in die
\E{unteren Atemwege} gelangt, spricht man von einer
(Fremdkörper-) \T{Aspiration}.

Nicht nur feste Körper, sondern auch Flüssigkeiten können
Auslöser einer Verlegung oder Aspiration sein. Besonders
Blut, Magensaft, Erbrochenes oder Speisebrei stellen
diesbezüglich eine Gefahr dar.
Besonders gefährdet hinsichtlich einer Atemwegsverlegung sind
Kinder, ältere Personen, bewusstseinseingschränkte Personen,
Patienten mit Schluckstörungen oder Allergiker (z.\,B.
Anschwellen der Schleimhäute nach einem Bienenstich).
Man unterscheidet zwischen einer \E{milden} (der Patient
kann noch sprechen und husten) und einer
\E{schweren Atemwegsverlegung} (Sprechen und Husten nicht
möglich) \citep{Erc2005Sec2,Erc2005Sec6,Erc2005DeSec2,Erc2005DeSec6}.



}
{%
\begin{itemize}
  \item Verlegung durch Fremdkörper
  \item Obere und untere Atemwege möglich
  \item Aspiration: untere Atemwege
  \item Besonders gefährdet:
  \begin{itemize}
    \item Kinder oder alte Personen
    \item Bewusstseinsgetrübte/-lose Personen
    \item Patienten mit Schluckstörungen
    \item Allergiker
  \end{itemize}
  \item Schweregrad:
  \begin{itemize}
    \item Milde Verlegung: \EE{Kann Sprechen \& Husten}
    \item Schwere Verlegung: Sprechen oder
    effektiv husten nicht möglich
  \end{itemize}
\end{itemize}
}
{./images/hirtler-aass/Fremdkoerper.pdf}
{Atemwegsverlegung (Schema). Der linke Bissen steckt in der Luftröhre fest, der
andere befindet sich in der Speiseröhre.}
{Hirtler.}




\Fazit{Die (Un-) Fähigkeit zu sprechen und zu husten ist das
wesentliche Kriterium um die Schwere einer Atemwegsverlegung
einzuschätzen.}

\TextSlide
{\TxABCDE}
{~
\begin{ABCDEText}
  \item[\ABCDE{1}] Achte auf Essenreste, Spielsachen
  (Murmeln!) und andere potentielle Fremdkörper.
  %%%
  \item[\ABCDE{2}] Oft gestikuliert der Patient
  entsprechend. Bei fortgeschrittener Hypoxie kann der
  Patient {\RedFlag}\,zyanotisch sein.
  %%%
  \item[\ABCDE{3}] Infolge einer Hypoxie
  kann es zu einer {\RedFlag}\,\E{Bewusstseinstrübung} kommen.
  %%%
  \item[\ABCDE{4}] Das Leitsymptom ist Atemnot bzw.
  Fremdkörpergefühl, zusammen mit einer entsprechenden
  Anamnese.
  %%%
  \item[\ABCDE{A}] Der Atemweg ist teilweise oder ganz
  verlegt. Eventuell ist ein pfeifendes Atemgaeräusch
  (\I{Stridor}) zu hören. Bei einer hochgradigen
  oder kompletten Verlegung ist das Atemgeräusch
  abgeschwächt oder fehlt völlig. Der Fremdkörper ist oft
  nicht sichtbar, da er sich oft unterhalb des Rachenraums
  befindet. 
  %%%
  \item[\ABCDE{B}] Je nach Ausmaß der Verlegung kann die
  Sauerstoffsättigung erniedrigt sein. Der Patient hat
  eine schnelle und erschwerte Atmung (\I{Tachypnoe}).
  \begin{itemize}
  \item \T[Atemwegsverlegung!milde]{Milde Verlegung}: Bei
  der milden Verlegung kann der Patient noch sprechen und effektiv husten. Er klagt
  meist über Atemnot und ein Fremdkörpergefühl, oft versucht er zu husten
  oder zu würgen. Oft ist auch ein pfeifendes Atemgeräusch
  während der Ein- und/oder Austamung zu hören
  (in-/exspiratorischer \E{Stridor}).
  \item \T[Atemwegsverlegung!schwere]{Schwere Verlegung}:
  Bei einer schweren Verlegung der Atemwege ist der Patient
  {\RedFlag}\,\EE{unfähig zu reden oder effektiv zu husten}.
  Eventuell kann er nur durch Nic\-k\-en oder Gesten kommunizieren. In besonders
  schweren Fällen kann es zu einem
  {\RedFlag}\,\EE{Atemstillstand} und nach kurzer Zeit zu einer
  {\RedFlag}\,\EE{Bewusstlosigkeit}, gefolgt von einem \EE{Kreislaufstillstand} kommen.
\end{itemize}
  %%%
  \item[\ABCDE{C}] Der Patient ist tachykard. Oft liegt eine Hypertonie aufgrund des
  Stress vor. 
  %%%
%   \item[\ABCDE{D}]
%   %%%
%   \item[\ABCDE{E}] 
%   %%%
%   \item[\ABCDE{\ldots}] 
%   %%%
  \item[\ABCDE{=}] Bei einer schweren Atemwegsverlegung ist
  der Patient vital bedroht. Bei der milden Verlegung ist
  die Einschätzung von den Befunden abhängig.
\end{ABCDEText}
}
{
\begin{ABCDESlide}
  \item[\ABCDE{1}] Essensreste, Spielsachen, etc.
  %%%
  \item[\ABCDE{2}] Gestik, evtl. {\RedFlag}\,Zyanose
  %%%
  \item[\ABCDE{3}] Evtl. {\RedFlag}\,\E{Bewusstseinstrübung}
  %%%
  \item[\ABCDE{4}] Atemnot, Fremdkörpergefühl
  %%%
  \item[\ABCDE{A}] %
  \begin{itemize}
    \item Fremdkörper oft nicht sichtbar
    \item Evtl. pfeifendes Atemgeräusch (Stridor)
  \end{itemize}
  %%%
  \item[\ABCDE{B}] %
  \begin{itemize}
    \item Atemnot
    \item Milde Verlegung: \EE{Kann Sprechen \& Husten}
    \item Schwere Verlegung:
    \begin{itemize}
      \item {\RedFlag}\,Unfähigkeit zu Sprechen oder
      effektiv zu husten
      \item Evtl. {\RedFlag}\,{Atemstillstand}, {\RedFlag}\,Bewusstlosigkeit,
      Hypoxie bis Kreislaufstillstand möglich
    \end{itemize}
  \end{itemize}
  %%%
  \item[\ABCDE{C}] HF $\uparrow$, RR $\uparrow$
  %%%
%   \item[\ABCDE{D}] 
  %%%
  \item[\ABCDE{=}] Bei schwerer Verlegung vital bedroht,
  sonst je nach Befund
\end{ABCDESlide}
}
{}
{}
{}




\TextSlide{Spezielle Techniken: Schläge zwischen die
Schulterblätter und Heimlich-Manöver}
{%
\begin{itemize}
  \item \T[Schulterblätter!Schläge zwischen die]{Schläge
  zwischen die Schulterblätter}:
  Durch kräftige Schläge auf den Rücken zwischen die
  Schulterblätter kann  ein festsitzender Fremdkörper evtl.
  gelockert und ausgehustet werden.
  \item \T{Heimlich-Manöver}:
  \label{topic:heimlich-manoever}%
  \Lang{Syn.} \I{Heimlich-Handgriff}. Beim Heimlich-Manöver
  versucht man durch Kompression des Bauchraums den Fremdkörper
  durch den entstandenen Überdruck aus den Atemwegen zu
  entfernen.
  
  Der Helfer umfasst von hinten den Oberbauch des Patienten und
  bildet mit der einen Hand eine Faust und umschliesst diese
  mit der anderen Hand. Nun legt er die Faust unterhalb der
  Rippen und des Brustbeins an und zieht sie dann ruckartig
  gerade nach hinten zu sich.
  
  \Ca{Verletzungsgefahr!} Durch Anwendung des
  Heimlich-Handgriffes kann es zu schweren inneren Verletzungen
  kommen (Magen, Leber, \ldots), die Anwendung ist daher auf
  bedrohliche Situationen beschränkt.
\end{itemize}
}{%
\begin{itemize}
  \item Schläge zwischen Schulterblätter
  \item Heimlich-Manöver
  \begin{itemize}
    \item Kompression des Bauchraumes → Überdruck
    \item Ansatzpunkt unterhalb des Rippenbogens, ruckartige Bewegung
    \item \Ca{Verletzungsgefahr!}
  \end{itemize}
\end{itemize}
}
{}
{}
{}


\begin{PictureSeries}{}{Bilderserie: Atemwegsverlegung}
\PictureSeriesRowWide{2}
{./images/hirtler-aass/Fremdkoerper.pdf}
{Atemwegsverlegung (Schema). Der linke Bissen steckt in der Luftröhre fest, der
andere befindet sich in der Speiseröhre. \SourcePicture{Hirtler.}}
{./images/wikipedia-pd/Abdominal_thrusts3.jpg}
{Heimlich-Manöver \SourcePicture{U.\,S. Naval Hospital, Okinawa, Japanese intern Naoko Uehara teaches
proper technique for treating a choking victim with the assistance of Hospital
Corpsman Second Class Ross Francis. Amanda M. Woodhead, USMC / WmCo, PD}}
{empty}
{}
{empty}
{}
\end{PictureSeries}


% M %%%%%%%%%%%%%%%%%%%%%%%%%%%%%%%%%%%%%%%%%%%%%%%%%%% M %
% Mechanische Atemwegsverlegung
\ProcedurePrintDefault{MT17091C}
 

% \Text{\dsliterary}{%
\nocite{Wuttig2009}
% }{}{}{}{}


\section{Asthma bronchiale}
\label{topic:asthma-bronchiale}

\Layout{60}{\TxBeschreibung{}}
{%
\DefinitionInPara{Das \T{Asthma bronchiale} ist eine chronische 
Erkrankung mit Episoden von Anfällen mit massiver Atemnot oder
Atembehinderung durch Verengung der Bronchien.}
Im akuten Asthmaanfall kommt es zu einer Verengung der
Bronchien und vermehrter Schleimbildung, die \EE{Ausatmung}
wird massiv erschwert und der Patient hat das Gefühl nicht
genug Luft zu bekommen.
}{%
Wiederholte Anfälle von Atemnot durch
\begin{itemize}
  \item Verengte Bronchien
  \item Vermehrte Schleimbildung
  \item Behinderte Ausatmung
\end{itemize}
}
{./images/asthma-paramedic.jpg}
{Patientin mit einem akuten Asthmaanfall. Sie sitzt auf einer Treppe und stützt sich nach hinten mit
den Armen ab. Die Erstmaßnahmen bei vital bedrohten
Patienten wurden ergriffen: Situationsrechte Lagerung,
Sauerstoffgabe, Reanimationsbereitschaft, Notarztnachforderung und Monitoring.\label{fig:asthma-paramedic}
}
{Gabriel}

Asthma ist bei den meisten Patienten oft schon bekannt und
kann bzw. muss erfragt werden. Viele Patienten haben auch
einen oder mehrere Sprays oder Inhalationsgeräte zur Dauer-
oder Akutbehandlung. 
Besonders betroffen vom Asthma bronchiale sind Patienten
mit bekannten Allergien. Folglich kann alles, was bei dem
Patienten eine Allergie verursacht auch einen Asthmaanfall
auslösen (Pollen, Gräser, Katzen, Zigarettenrauch, \ldots). 

Ist bei Eintreffen des Rettungsdienstes noch keine
Besserung der Atemnot eingetreten, ist diese Atemnot als
massiv und lebensbedrohend einzuschätzen. Es ist nicht
damit zu rechnen, dass der Zustand von alleine wieder
vergeht und bedarf einer medikamentösen Behandlung durch einen
Notarzt. 

\TextSlide{Status Asthmaticus}{%
\index{Status!asthmaticus}
\label{topic:status-asthmaticus}
Der
\I[Status!asthmaticus]{Status asthamticus} ist eine gefürchtete Komplikation des
Asthmaanfalles bei der die Atemnot nicht wieder vergeht. Es
besteht die Gefahr, dass der Patient nach einiger Zeit keine
Kraft mehr zum Atmen hat (Erschöpfung). Der Status
asthmaticus kann tödlich enden.
}{
\begin{itemize}
  \item Lebensgefährliche Komplikation
  \item Anfall hört nicht mehr auf
\end{itemize}
}{}{}{}




\TextSlide
{\TxABCDE}
{~
\begin{ABCDEText}
  \item[\ABCDE{1}] Kalte Temperaturen oder Allergene
  (Gräser, Pollen, \ldots) begünstigen einen Asthma-Anfall.
  %%%
  \item[\ABCDE{2}] Der Patient nimmt meist automatisch und unbewusst eine
  Position ein, die seine Atemhilfsmuskulatur unterstützt: Er
  sitzt aufrecht und stützt sich oft nach vorne oder hinten
  ab.
  %%%
  \item[\ABCDE{3}] Im schweren Anfall kann es zu
  {\RedFlag}\,\E{Bewusstseinsstörungen} kommen.
  %%%
  \item[\ABCDE{4}] Das Leitsymptom ist die Atemnot.
  %   %%%
  %   \item[\ABCDE{A}]
  %%%
  \item[\ABCDE{B}] Die Atemfrequenz ist erhöht
  (\E{Tachypnoe}), die {\RedFlag}\,\E{Sauerstoffsättigung}
  kann erniedrigt sein. Bei einer manifesten Hypoxie kann eine
  {\RedFlag}\,\E{Zyanose} beobachtbar sein.
  
  Oft ist ein pfeifendes,
  brummendes oder giemendes Atemgeräusch, besonders während
  der Ausatmung (\E{Exspiration}, \E{exspiratorischer
  Stridor}) wahrnehmbar. Die Exspiration ist ausserdem
  verlängert. Trockener Husten ist häufig.
  
%   \Z{Unterschätze nie die Schwere eines Asthmaanfalles!
%   Gefährlich sind besonders die \protect\enquote*{ruhigen},
%   wenig klagenden Patienten, diese haben oft ein anderes
%   Luftnot-Empfinden. Diese Patienten sind Hochrisiko-Patienten
%   und haben auch mit der \ce{O2}-Gabe Probleme! \ce{O2}-Gabe nur
%   bis max. 95\PC* Sauerstofsättigung!
%   \citep{Olschewski2009}\A{\footnote{Vortrag Univ-Prof.Dr. Horst Olschewski, 5. AGN-Kongress 2009; Notfall Rettungsmed (2009) 12:322}}}
  
  Im \EE{schweren Anfall} ist oft \E{kein}, oder nur mehr ein
  \E{leises Atemgeräusch} hörbar, da kaum mehr Luft bewegt
  wird \Z{(\I{silent chest})}. Die
  {\RedFlag}\,Sauerstoffsättigung ist stark vermindert und der
  Patient ist {\RedFlag}\,\E{zyanotisch}. Im Endstadium kann der 
  Patient zu {\RedFlag}\,erschöpft für eine ausreichende
  \E{Atemarbeit} sein. 
  
  Bei der Einschätzung der Atmung ist es wichtig auf die
  \EE{Zeichen einer erschwerten Atemarbeit} zu
  achten:
  \begin{itemize}
    \item Einsatz der Atemhilfsmuskulatur
    \item Einziehungen an den Rippen
    \item Erschöpfungszeichen
  \end{itemize}
  Es besteht die Gefahr der {\RedFlag}\,\EE{Erschöpfung der
  Atemmuskulatur} und einer reduzierten Atemarbeit.
  %%%
  \item[\ABCDE{C}] Die Herzfrequenz ist meist erhöht
  (\E{Tachykardie}).
%   
%   
%   %%%
%   \item[\ABCDE{D}]
%   %%%
%   \item[\ABCDE{E}]
%   %%%
%   \item[\ABCDE{\ldots}]
  %%%
  \item[\ABCDE{=}] Bei schwerer Atemnot besteht 
  eine vitale Bedrohung. Bei Erschöpfung der Atemarbeit
  besteht akute Lebensgefahr!
\end{ABCDEText}
}
{
\begin{ABCDESlide}
  \item[\ABCDE{1}] Kälte, Allergene
  %%%
  \item[\ABCDE{2}] Einsatz der Atemhilfsmuskulatur
  %%%
  \item[\ABCDE{3}] Schwerer Anfall:
  {\RedFlag}\,Bewusstseinsstörungen
  %%%
  \item[\ABCDE{4}] Akute Atemnot
%   %%%
%   \item[\ABCDE{A}] 
  %%%
  \item[\ABCDE{B}] Akute Atemnot mit trockenem
  Husten, Verlängerte Ausatmungsphase, AF $\uparrow$,
  Exspiratorisches Atemgeräusch (Stridor)
  
  Schwerer Anfall: Kein oder leises Atemgeräusch,
  {\RedFlag}\,Sp\ce{O2} $\downarrow\downarrow$
  
  {Zeichen einer erschwerten Atemarbeit}:
  \begin{itemize}
    \item Einsatz der Atemhilfsmuskulatur
    \item Einziehungen an den Rippen
    \item Erschöpfungszeichen
  \end{itemize}
 
  Endstadium: {\RedFlag}\,Atemarbeit erschöpft
  %%%
  \item[\ABCDE{C}] HF $\uparrow$
  
  Schwerer Anfall: HF $\downarrow$
  %%%
  \item[\ABCDE{D}] 
  Schwerer Anfall: RR $\downarrow$
%   %%%
%   \item[\ABCDE{E}] 
%   %%%
%   \item[\ABCDE{\ldots}] 
  %%%
  \item[\ABCDE{=}] Schwere Atemnot: Vital bedroht.
  Erschöpfung: Atemarbeit $\downarrow$ $\rightarrow$ Akute Lebensgefahr
\end{ABCDESlide}
}
{}
{}
{}




\TextSlide
{{\TxSAMPLER}}
{%
\begin{SAMPLERText}
  \item[\SAMPLER{A}] Patienten mit allergischem Asthma haben
  häufig Allergien gegen Gräser, Pollen, Tierhaare usw.
  \item[\SAMPLER{M}] Oft werden Sprays bzw. Inhalatoren zur
  Dauertherpaie und für den Anfall eingesetzt. Manche
  Medikamente, insbesonders Schmerzmedikamente, können
  Anfälle auslösen.%
  \ZFootnote{Besonders nichtsteroidale Antirheumatika (ASS,
  Aspirin{\texttrademark}, etc.) und Betablocker
  (Antiarrhythmika, Blutdruckmittel) können Asthma-Anfälle
  auslösen.} 
  Nahm der Patient in letzter Zeit (neue)
  Medikamente?
  \item[\SAMPLER{P}] Meist ist Asthma vorbekannt. 
  \item[\SAMPLER{L}] Letzte Spray-Einnahme erfragen!
  \item[\SAMPLER{E}] Häufige Auslösefaktoren sind Stress,
  Anstrengung und Allergenexposition.
%   \item[\SAMPLER{R}]
\end{SAMPLERText}
}
{
\begin{SAMPLERSlide}
%   \item[\SAMPLER{S}]
  \item[\SAMPLER{A}] Gräser, Pollen, Tierhaare, \ldots
  \item[\SAMPLER{M}] Sprays, Inhalatoren. (Neue) Medikamente in letzter Zeit?
  \item[\SAMPLER{P}] Asthma oft vorbekannt
  \item[\SAMPLER{L}] Spray-Einnahme?
  \item[\SAMPLER{E}] Stress, Anstrengung, Allergene
%   \item[\SAMPLER{R}]
\end{SAMPLERSlide}
}
{}
{}
{}




% \begin{figure}[h]
%     \includegraphics{./images/copd-bronchus.png}
%     
%     \Captionof{figure}{Bronchus bei COPD.
%     \AuthNotesPicLegalNotOk{copd-bronchus}{}}
% \end{figure}


% M %%%%%%%%%%%%%%%%%%%%%%%%%%%%%%%%%%%%%%%%%%%%%%%%%%% M %
% Akuter Asthmaanfall
\ProcedurePrintDefault{MJ46XX1C}


% \clearpage % TEMP temp clearpage
\section{Chronische
Bronchitis und COPD}
\label{topic:copd}
\nocite{Renz-Polster2006} 




\TextSlide{\TxBeschreibung: COPD}{%
\DefinitionInPara{Die \T{chronische Bronchitis} ist
eine chronische entzündliche Schleimhautschädigung  der unteren Atemwege.}
\DefinitionInPara{Die \T{COPD}
(\T[Lungenerkraknkung!chronisch-obstruktive]{Chronische Obstruktive
Lungenerkrankung}\ZFootnote{\Lang{engl.} Chronic
obstructive pulmonary disease\index{Chronic
obstructive pulmonary disease|Z})} ist eine chronische entzündliche
Schleimhautschädigung, welche eine zunehmende obstruktive 
Atemwegseinschränkung aufweist.}
Die \Abk{COPD} ist durch eine voranschreitende
Verschlechterung der Atemleistung gekennzeichnet.
Am Anfang steht die chronische Bronchitis, welche durch
Husten mit schleimigem Auswurf gekennzeichnet ist
(\enquote*{\I{Raucherhusten}}). Es kommt dabei zu einer gesteigerten
Entzündungsantwort auf eingeatmete Stoffe (Zigarettenrauch,
Umweltschadstoffe, \ldots). Wenn man in der
Lungenfunktionsuntersuchung eine Atmungsseinschränkung
nachweisen kann, spricht man von der chronisch-obstruktiven
Lungenerkrankung (\Abk{COPD}). 
Diese geht mit bleibenden Veränderungen der unteren Atemwege einher, 
diese sind in der Grafik  \prettyref{fig:copd-veraenderungen} dargestellt.

\Z{Durch die erschwerte Ausatmung kommt es zu einer
chronischen Überblähung der Lungenbläschen und zu einer \protect\enquote{Faß-förmigen}
Verformung des Brustkorbes. Im Endstadium zeigen sich
Zeichen einer Rechtsherzinsuffizienz
(\prettyref{topic:rechtsherzinsuffizienz}) aufgrund einer
Störung im Lungenkreislauf.}

\subparagraph{Exazerbationen}
Kommt zu der ohnehin schweren Grunderkrankung noch ein
erschwerender Faktor hinzu, z.\,B. eine Infektion der
Atemwege, kann es zu einer plötzlichen Verschlechterung
kommen, zur \E{Exazerbation}. Diese ist meist durch
vermehrte Atemnot, Husten und Auswurf gekennzeichnet.
}{
\begin{itemize}
  \item Husten mit Auswurf, anfänglich
  \protect\enquote*{Raucherhusten}
  \item Voranschreitende, sich verschlechternde \E{Ateminsuffizienz}:
  \begin{compactitem}
    \item Belastungsdyspnoe (bei schwererer \Abk{COPD}
    auch in Ruhe)
    \item Ausatmung (Exspiration) erschwert
    \item Brummendes, pfeifendes od. giemendes Ausatemgeräusch
    \item Evtl. chronische Zyanose
    \item Evtl. Heimsauerstoff
    \item Sp\ce{O2} chronisch niedrig
  \end{compactitem}
  \item Exazerbation:
\begin{itemize}
  \item Plötzliche Verschlechterungen
  \item Meist Infekt-bedingt
  \item Vermehrte Atemnot,
  Husten und Auswurf
\end{itemize}
\end{itemize}
}{}{}{}



\TextSlide{Probleme mit Sauerstoff bei
COPD-Patienten}
{
\index{Sauerstoff!COPD-Patient}
Normalerweise wird der Atemantrieb über die
\ce{CO2}-Konzentration im Blut gesteuert. Die
\ce{CO2}-Konzentration ist bei \Abk{COPD}-Patienten jedoch chronisch
hoch, der Körper gewöhnt sich an diesen Zustand. Er
regelt dann den Atemantrieb direkt über den \ce{O2}-Gehalt im Blut.
Wird der \ce{O2}-Gehalt plötzlich künstlich erhöht, fehlt
der \index{COPD!Atemantrieb}Atemantrieb und der Patient kann
\EE{Bewusstseinsstörungen} (\I{\ce{CO2}-Narkose}) und
im Extremfall einen \EE{Atemstillstand} erleiden.

}{
\Ca{Achtung!\\
Bewusstseinsstörungen bei übermäßiger
Sauerstoffgabe möglich!}

\SymbolText
}{}{}{}



\begin{figure}[H]

\includegraphics[width=\linewidth]{./images/hirtler-aass/Alveolen-Bronchus-COPD-edited2.png}
\Captionof{figure}{\label{fig:copd-veraenderungen}Veränderung der
Atemwege bei der COPD. \TabSub{Links:} Schema eines gesunden
Bronchus und einer gesunden Alveole. \TabSub{Oben rechts:} Bei der COPD sind
die kleinen Luftwege verschleimt und verengt. \TabSub{Unten rechts:} Die
Lungenbläschen (Alveolen) sind überbläht, weil die Luft nur erschwert wieder entweichen kann.
\SourcePicture{Hirtler.}}
\end{figure}






% 1-E
\TextSlide
{\TxABCDE: COPD}
{~
\begin{ABCDEText}
  \item[\ABCDE{1}] Evtl. kalte Umgebung (kalte Luft führt
  zur Verengung der Bronchien)
  \item[\ABCDE{2}] Einsatz der Atemhilfsmuskulatur, Mühe
  beim Atmen, evtl. Heimsauerstoff, oft wirkt der Patient
  ängstlich.
  \item[\ABCDE{3}] Bei schweren Anfällen kann es aufgrund
  der Hypoxie zu {\RedFlag}\,\EE{Bewusstseinsstörungen}
  kommen.
  \item[\ABCDE{4}] Atemnot
  \item[\ABCDE{A}] Siehe \ABCDE{B}
  \item[\ABCDE{B}] Je nach Schweregrad kommt es zu Zeichen
  einer Atemwegsverlegung der \E{unteren}
  Atemwege  (Obstruktion, durch Verengung der Bronchien und
  Schleimproduktion) und Ateminsuffizienz:
  Die \E{Ausatmung} (Exspiration) ist erschwert, der Patient klagt
  über \E{Belastungsdyspnoe}, ein Engegefühl, \E{Husten} und
  hat Tachypnoe und evtl. {\RedFlag}\,\EE{Zyanose}. Man kann
  oft ein
  \E{brummendes, pfeifendes oder giemendes Atemgeräusch}
  hören.
  
  Die \EE{Sauerstoffsättigung} bei einem COPD-Patienten ist
  auch ohne akute Symptome meist \E{deutlich vermindert}
  ($\sim$\,{90}\PC*). Der sonst übliche Normalwert
  gilt hier nicht! Bei fortgeschrittener Ateminsuffizienz
  bekommt der Patient oft eine \E{Heimsauerstofftherapie}
  verordnet. Bei übermäßiger Sauerstoffgabe kann es zu
  {\RedFlag}\,\EE{Bewusstseinsstörungen kommen}
  (\I{\ce{CO2}-Narkose})!
  \item[\ABCDE{C}] Evtl. tachkard und hyperton
%   \item[\ABCDE{D}]
%   \item[\ABCDE{E}]
  \item[\ABCDE{=}] Bei {\RedFlag}\,schwerer Atemnot
  oder {\RedFlag}\,Bewusstseinsstörungen vitale Bedrohung.
  \item[\ABCDE{\ldots}] Als Nebenbefund können Infektzeichen
  vorliegen (erhöhte Körpertemperatur).
\end{ABCDEText}
}
{
\begin{ABCDESlide}
  \item[\ABCDE{1}] Kalte Umgebung
  \item[\ABCDE{2}] Atemhilfsmuskulatur, Anstrengung beim
  Atmen, ängstlich
  \item[\ABCDE{3}] Evtl. {\RedFlag}\,Bewusstseinsstörungen
  \item[\ABCDE{4}] Atemnot
  \item[\ABCDE{A}] Siehe \ABCDE{B}
  \item[\ABCDE{B}] Atemnot! Ausatmung erschwert, Husten
  Tachypnoe, {\RedFlag}\,Zyanose;
  
  Brummendes, pfeifendes Atemgeräusch
  
  Sp\ce{O2} chronisch niedrig ($\sim$\,{90}\PC*): Normalwert
  gilt nicht!
  \item[\ABCDE{C}] HF $\uparrow$, RR $\uparrow$
%   \item[\ABCDE{D}] 
%   \item[\ABCDE{E}] 
  \item[\ABCDE{=}] Vitale Bedrohung bei schwerer Atemnot und 
  Bewusstseinsstörungen
  \item[\ABCDE{\ldots}] Evtl. erhöhte Körpertemperatur
\end{ABCDESlide}
}
{}
{}
{}

\TextSlide
{{\TxSAMPLER}: COPD}
{
\begin{SAMPLERText}
%   \item[\SAMPLER{S}] 
%   \item[\SAMPLER{A}]
  \item[\SAMPLER{M}] Inhalatoren bzw. Sprays zur Dauer- und
  Akuttherapie, Heimsauerstoff. \Z{Häufig Kortisonpräparate;
  bei vorbekanntem Infekt wurden oft schon Antibiotika
  verschrieben.}
  \item[\SAMPLER{P}] Eine COPD ist normalerweise vorbekannt.
  \item[\SAMPLER{L}] Letzte Spray-Einnahme?
%   \item[\SAMPLER{E}] 
%   \item[\SAMPLER{R}]
\end{SAMPLERText}
}
{
\begin{SAMPLERSlide}
%   \item[\SAMPLER{S}]
%   \item[\SAMPLER{A}]
  \item[\SAMPLER{M}] Inhalatoren, evtl. Heimsauerstoff.
  \item[\SAMPLER{P}] COPD vorbekannt
  \item[\SAMPLER{L}] Letzte Spray-Einnahme?
%   \item[\SAMPLER{E}] 
%   \item[\SAMPLER{R}]
\end{SAMPLERSlide}
}
{}
{}
{}



\begin{figure}[h]
\begin{center}
% \includegraphics{./images/copd-stufen.png}
\includegraphics[width=\linewidth]{./images/hirtler-aass/COPD-Stufen_edited.pdf}

\Captionof{figure}{Eine COPD
entsteht nicht plötzlich: Ein COPD{\textquotesingle}ler hat eine
\protect\enquote{Karriere} hinter sich. \SourcePicture{Hirtler}}
\end{center}
\end{figure}


% M %%%%%%%%%%%%%%%%%%%%%%%%%%%%%%%%%%%%%%%%%%%%%%%%%%% M %
% COPD-Exazerbation
\ProcedurePrintDefault{MJ44190C}


\nocite{Knacke200601}


\section{Lungenembolie}
\label{topic:lungenembolie}



\Layout{60}{{\TxBeschreibung} und Ursachen}
{%
\DefinitionInPara{Als \T{Lungenembolie} (\I{Pulmonalembolie}) wird
der \EE{Verschluss einer Lungenarterie}, oft  durch Einschwemmen eines 
\EE{Gerinnsels} (\I{Embolus}), bezeichnet.}
Der Embolus, der die Lungenarterien verstopft, entsteht
normalerweise nicht in den Lungengefäßen selbst, sondern er
wird über die Blutbahn angeschwemmt. Sehr oft löst sich ein Thrombus bei
einer \EE{tiefen Beinvenenthrombose} und schwimmt im Blut bis in
die Lunge. Daher gelten die
Risikofaktoren der tiefen Beinvenenthrombose
(Immobilisation, lange Reisen, \ldots) auch bei der
Lungenembolie. 
Auch bei \EE{Vorhofflimmern} können sich Gerinnsel im 
rechten Vorhof des Herzens bilden und eine Embolie verursachen.
Ebenso können
Gerinnungstörungen oder Krebserkrankungen zur Bildung von 
Blutgerinnseln führen.
}{
\begin{itemize}
  \item Verschluss einer Lungenarterie
  \item Oft eingeschwemmt, vor allem:
  \begin{itemize}
    \item Tiefe Beinvenenthrombose
    \item Rechter Vorhof bei Vorhofflimmern
  \end{itemize}
\end{itemize}
}
{./images/hirtler-aass/Thrombus-Embolie.pdf}
{Herkunft des Thrombus.}
{Hirtler.}




\TextSlide
{\TxABCDE: Lungenembolie}
{~
\begin{ABCDEText}
  \item[\ABCDE{1}] Lässt die Szenarie Hinweise auf eine Bettlägrigkeit
  (Pflegebett, \ldots) erkennen? Gibt es Hinweise auf eine Fernreise (herumstehende Koffer, {\ldots})?
  \item[\ABCDE{2}] Kann man beim Patienten Hinweise auf Bettlägrigkeit (Alter,
  Gipsverband, körperlicher Zustand, \ldots) erkennen?
%   \item[\ABCDE{3}]
  \item[\ABCDE{4}] Das Leitsymptom ist die
  {\RedFlag}\,\EE{Atemnot}, eventuell in Kombination mit
  (atemabhängigen) {\RedFlag}\,Thorax- oder Rückenschmerzen. 
  
  Eine Unterscheidung zwischen einer Lungenembolie und
  einem Akuten Koronarsyndrom ist oft nicht möglich. 
%   \item[\ABCDE{A}]
  \item[\ABCDE{B}] Meist atmet der Patient schnell (Tachypnoe) und klagt über
  einen \EE{atemabhängigen {\RedFlag}\,Thoraxschmerz}.
  \item[\ABCDE{C}]  Es kommt zu
  einer \E{Tachykardie} und Angstgefühlen, oft findet man
  deutlich hervortretende, \EE{gestaute Venen}
  am Hals aufgrund des Blutrückstaus. Evtl. kann man auch einen
  \E{arrhythmischen Puls} aufgrund eines Vorhofflimmerns, ertasten.
  Der Blutdruck kann normal oder vermindert sein. Bei einer
  schweren Lungenembolie kann es zu einem
  {\RedFlag}\,\EE{kardiogenen Schock} kommen, es ist daher
  wichtig auf {\RedFlag}\,Schockzeichen zu achten!
%   \item[\ABCDE{D}]
%   \item[\ABCDE{E}]
  \item[\ABCDE{=}] Bei {\RedFlag}\,anhaltender Atemnot oder
  {\RedFlag}\,Thoraxschmerzen besteht eine vitale Bedrohung.
  \item[\ABCDE{\ldots}] Je nach Ursache der Embolie und
  Herkunft des Embolus kann
  man Zeichen anderer Krankheitsbilder wahrnehmen: Bei einer
  \EE{tiefen Beinvenenthrombose} ist ein Bein geschwollen,
  überwärmt, rot/rot-bläulich verfärbt und schmerzhaft
  (\prettyref{topic:thrombose}).
  
\end{ABCDEText}
}
{
\begin{ABCDESlide}
  \item[\ABCDE{1}] Bettlägrigkeit? Fernreise?
  \item[\ABCDE{2}] Bettlägrigkeit?
%   \item[\ABCDE{3}] 
  \item[\ABCDE{4}] Dyspnoe, (atemabhängige)
  {\RedFlag}\,Thorax- od. Rückenschmerzen, Angstgefühl
  
  Unterscheidung zu ACS oft nicht möglich
%   \item[\ABCDE{A}] 
  \item[\ABCDE{B}] {\RedFlag}\,Atemnot, Tachypnoe
  \item[\ABCDE{C}] HF $\uparrow$, evtl. arrhythmischer Puls
  (Vorhofflimmern), RR \SymbolFindingNormal od.
  {\RedFlag}\,$\downarrow$

  Gestaute Halsvenen (Blutrückstau), evtl. arrhythmischer Puls bei Vorhofflimmern
  
  {\RedFlag}\,Schockzeichen?
% 
%   \item[\ABCDE{D}] 
%   \item[\ABCDE{E}] 
  \item[\ABCDE{=}] Bei bestehender Atemnot oder
  Thoraxschmerzen vitale Bedrohung.
  \item[\ABCDE{\ldots}] Venöser Verschluss?
\end{ABCDESlide}
}
{}
{}
{}

\TextSlide
{{\TxSAMPLER}: Lungenembolie}
{
\begin{SAMPLERText}
%   \item[\SAMPLER{S}]
%   \item[\SAMPLER{A}]
  \item[\SAMPLER{M}] Bettlägrigen Patienten werden zur
  Thromboseprophylaxe oft gerinnungshemmende Medikamente
  verschrieben (\I{Thrombosespritzen}). Wurden diese
  genommen?
  \item[\SAMPLER{P}]
  Oft ist bei den Patienten \EE{Vorhofflimmern} als
  chronische Erkrankung bekannt. Patienten, welche bereits
  schon einmal eine Thrombose oder eine Embolie hatten,
  haben ein erhöhtes Risiko neuerlich daran zu erkranken. Im
  Rahmen von Krebserkrankungen kann es gehäuft zu Thrombosen
  und Embolien kommen.
  % \item[\SAMPLER{L}]
  \item[\SAMPLER{E}] Fernreise? Langes Sitzen/Liegen? Wenig
  getrunken?
  \item[\SAMPLER{R}] \EE{Immobilisation}, z.\,B. durch Bettlägrigkeit,
  Gipsverbände oder lange Reisen, innerhalb der letzten 4
  Wochen, erhöht das Risiko einer Thrombose und damit auch das
  einer Lungenembolie. Ebenso erhöht eine bereits durchgemachte 
  Lungemembolie oder Thrombose das Risiko. 
  \Z{Weiters tritt die Lungenembolie gehäuft in der Schwangerschaft
  auf. Rauchen sowie die Einnahme der \protect\enquote{Pille} erhöhen
  ebenso das Risiko.}
\end{SAMPLERText}
}
{
\begin{SAMPLERSlide}
%   \item[\SAMPLER{S}]
%   \item[\SAMPLER{A}]
  \item[\SAMPLER{M}] Thrombosespritzen
  \item[\SAMPLER{P}] Vorhofflimmern, Thrombose, Embolie,
  Krebserkrankungen
%   \item[\SAMPLER{L}]
  \item[\SAMPLER{E}] Fernreise? Langes Sitzen/Liegen? Wenig
  getrunken?
  \item[\SAMPLER{R}] Immobilisation, Vorhofflimmern, vorangegange Thrombose oder Embolie, 
  \Z{Schwangerschaft}, \Z{Rauchen}, \Z{Verhütungspille}
\end{SAMPLERSlide}
}
{}
{}
{}





% M %%%%%%%%%%%%%%%%%%%%%%%%%%%%%%%%%%%%%%%%%%%%%%%%%%% M %
% Lungenembolie
\ProcedurePrintDefault{MI26XX0C}




\section{Lungenödem}
\label{topic:lungenoedem}

\nopagebreak

\TextSlide{{\TxBeschreibung} und Ursachen}
{%
\DefinitionInPara{Als \T{Lungenödem} wird die 
Flüssigkeitsansammlung im Zwischengewebsraum und in den
Alveolen und der Lungen bezeichnet.}
Ein Lungenödem kann entweder aufgrund einer
\EE{Herzschwäche} (\E{kardial}) oder aufgrund
  von \EE{Inhalation giftiger Substanzen} (\E{toxisch})
entstehen:
\begin{itemize}
  \item Beim \EE{kardialen Lungenödem} versagt die Auswurfleistung
  des linken Herzens (\EE{Linksherzinsuffizienz}), dadurch
  kommt es zu einem \EE{Rückstau} des Blutes in die Lunge.
  Aufgrund der Druckerhöhung tritt Flüssigkeit in das Gewebe
  über, es bilden sich Ödeme im Gewebe bzw.
  \EE{Flüssigkeitsansammlungen} in den Alveolen. Auslöser kann
  zum Beispiel ein Herzinfarkt sein, der die Funktion des
  Hermuskels massiv in Mitleidenschaft zieht.
  \item Beim \EE{toxischen  Lungenödem} kommt es aufgrund reizender
  Stoffe zu einer \EE{Entzündung}, und damit zu einem
  \EE{Flüssigkeitsaustritt} im Lungengewebe.
\end{itemize}
In beiden Fällen behindert die Flüssigkeit den
\EE{Gasaustausch} massiv, beim kardialen Lungenödem kommt
noch das begleitende Herzproblem dazu.

}{
\begin{itemize}
  \item Flüssigkeitsansammlung im Zwischengewebsraum und in 
  Alveolen der Lunge
  \item Arten
  \begin{itemize}
    \item \EE{Kardial}: (Links-)Herzinsuffizienz (Blutrückstau); 
    z.\,B. bei MCI, zusätzlicher Belastung bei vorgeschädigtem Herz
    (Infekt), \ldots
    \item \EE{Toxisch}: Inhalation giftiger Substanzen
  \end{itemize}
  \item Gasaustausch beeinträchtigt
\end{itemize}

}{}{}{}




\TextSlide
{\TxABCDE: Lungenödem}
{~
\begin{ABCDEText}
  \item[\ABCDE{1}] Toxisches Lungenödem: Besonders auf
  Umgebungssicherheit achten: Offene Lacke,
  Lösungsmittel, Dämpfe, Chemikalien, Warnschilder, \ldots.
  \item[\ABCDE{2}] Der Patient sitzt, oft sind die Pölster
  im Bett aufgetürmt.
%   \item[\ABCDE{3}] 
  \item[\ABCDE{4}] Atemnot
  \item[\ABCDE{A}] Es kann zu schaumigem Auswurf kommen.
  \item[\ABCDE{B}] \index{Brodelnde Atemgeräusche}%
  Das Leitsymptom ist das
  {\RedFlag}\,\EE{brodelnde Atemgeräusch}, welches in
  schweren Fällen bereits ohne Hilfsmittel hörbar und
  typisch für das Lungenödem ist. Es klingt, als würde jemand
  mit einem Strohhalm in ein Glas Wasser hineinblasen. Im
  Extremfall kann man es bereits \E{auf mehrere Meter
  Entfernung} wahrnehmen.
  Weitere häufige Symptome sind anhaltende {\RedFlag}\,\EE{Atemnot} und
  {\RedFlag}\,\EE{Zyanose}.
  Der Patient bekommt oft nur \E{aufrecht sitzend} genügend
  Luft \Z{(Orthopnoe)}, im Liegen verschlechtert sich die
  Atemnot deutlich.
  \item[\ABCDE{C}] Die Herzfrequenz und der Blutdruck können
  in jede Richtung verändert sein (Herzrhythmusstörungen
  oder ein hypertensiver Notfall können ein kardiales
  Lungenödem auslösen!). Im Verlauf kann es zu einem
  {\RedFlag}\,\EE{kardiogenen Schock} kommen.
%   \item[\ABCDE{D}]
%   \item[\ABCDE{E}] 
  \item[\ABCDE{=}] Patienten mit einem {\RedFlag}\,deutlich
  hörbaren Lungenödem (v.\,a. bei kardialer Ursache) sind schwer und \EE{lebensbedrohlich
  krank}. Das Herz verbraucht oft schon die allerletzten Reserven. In
  schweren Fällen kann schon die Aufregung durch den Transport
  im Stiegenhaus zum Auto zu einem Herz-Kreislaufstillstand
  führen.
  
  Es kommt tatsächlich relativ häufig vor, dass diese
  Patienten \EE{im Stiegenhaus} oder \EE{im Fahrzeug}
  \EE{reanimationspflichtig} werden. \Ca{Solche Patienten müssen
  daher \EE{an Ort und Stelle} eine entsprechende (ärztliche)
  Therapie erhalten bevor sie transportiert werden können.}
  
  Bei leichten Formen des Lungenödems sind die
  Atemgeräusche nicht mit freiem Ohr, sondern
  nur mittels Stethoskop zu hören.  Die Patienten zeigen auch
  sonst keine Alarmsymptome, sondern nur Ödeme und eine mäßige
  Kurzatmigkeit. Diese Patienten sind i.\,d.\,R. nicht
  unmittelbar vital bedroht, ein bestmögliches Monitoring ist jedoch wichtig.
  \item[\ABCDE{\ldots}] Beidseitige Beinödeme sind typisch.
\end{ABCDEText}
}
{
\begin{ABCDESlide}
  \item[\ABCDE{1}] Umgebung sicher? Chemikalien?
  Warnschilder?
  \item[\ABCDE{2}] Sitzend, aufgetürmte Pölster
%   \item[\ABCDE{3}]
  \item[\ABCDE{4}] Atemnot
  \item[\ABCDE{A}] Evtl. schaumiger Auswurf
  \item[\ABCDE{B}] {\RedFlag}\,Brodelndes Atemgeräusch

  Atemnot (besonders im Liegen), evtl. {\RedFlag}\,Zyanose
  
  Atemnot im Liegen schlechter
  \item[\ABCDE{C}] HF $\uparrow\downarrow$, RR
  $\uparrow\downarrow$
  
  Evtl. kardiogener Schock
%   \item[\ABCDE{D}]
%   \item[\ABCDE{E}]
  \item[\ABCDE{=}]
  \begin{itemize}
    \item {\RedFlag}\,Hörbares Lungenödem:
    \Ca{Lebensgefahr!}
    \item Leichte Form: Keine Alarmsymptome
  \end{itemize}
  \item[\ABCDE{\ldots}] Beidseitige Beinödeme
\end{ABCDESlide}
}
{}
{}
{}

\TextSlide
{{\TxSAMPLER}: Lungenödem}
{
\begin{SAMPLERText}
%   \item[\SAMPLER{S}] 
%   \item[\SAMPLER{A}]
  \item[\SAMPLER{M}] Herzpatienten nehmen oft zahlreiche
  Medikamente ein (Blutdruck, Entwässerung, Blutfette,
  \ldots).
  \item[\SAMPLER{P}] Häufig berichten die Patienten über
  Herz-Kreislauferkrankungen wie z.\,B. Herzinsuffizienz,
  Z.\,n. Herzinfarkt, Bluthochdruck oder
  Herzklappenerkrankungen.
%   \item[\SAMPLER{L}]
  \item[\SAMPLER{E}] War der Patient chemischen Substanzen ausgesetzt?
  Ertrinkungsunfall?
  
  Erschwerender Faktor bei vorbestehender Herzerkrankung
  (Infekt, \ldots)?
%   \item[\SAMPLER{R}]
\end{SAMPLERText}
}
{
\begin{SAMPLERSlide}
%   \item[\SAMPLER{S}]
%   \item[\SAMPLER{A}]
  \item[\SAMPLER{M}] Zahlreich
  \item[\SAMPLER{P}] Herz-Kreislauferkrankungen: Herzinsuffizienz,
  Z.\,n. Herzinfarkt, Bluthochdruck,
  Herzklappenerkrankungen, \ldots
%   \item[\SAMPLER{L}]
  \item[\SAMPLER{E}] Exposition? Ertrinkungsunfall?
  Erschwerender Faktor beim vorbestehender Herzerkrankung
  (Infekt, \ldots)?
%   \item[\SAMPLER{R}]
\end{SAMPLERSlide}
}
{}
{}
{}



% M %%%%%%%%%%%%%%%%%%%%%%%%%%%%%%%%%%%%%%%%%%%%%%%%%%% M %
% Akutes Lungenödem, leicht
\ProcedurePrintDefault{MJ81XX1C}


% M %%%%%%%%%%%%%%%%%%%%%%%%%%%%%%%%%%%%%%%%%%%%%%%%%%% M %
% Akutes Lungenödem, schwer
\ProcedurePrintDefault{MJ81XX2C}



\section{Hyperventilationssyndrom und -tetanie}
\label{topic:hyperventilationstetanie}
\label{topic:hyperventailationssyndrom}
\index{Hyperventilation!-stetanie}
\index{Hyperventilation!-ssyndrom}
\index{Hyperventilation!psychisch bedingte --}



\TextSlide{{\TxBeschreibung} und Ursachen}{%
\DefinitionInPara{Als \T{Hyperventilationssyndrom} wird die 
übermäßig tiefe und schnelle Atmung \E{ohne körperliche Ursache} 
bezeichnet.}
Durch die Atmung wird der Säure-Basen-Haushalt des Körpers
wesentlich beeinflusst
(\prettyref{topic:saeure-basen-haushalt}). Atmet ein sonst
gesunder Mensch übermäßig tief und schnell, so wird sein
Blut-pH sehr rasch basisch. Dadurch geraten die
\E{Elektrolyte} im Blut in ein Ungleichgewicht, es kommt zu
Symptomen.

Das Hyperventilationssyndrom ist meistens psychisch bedingt
und Folge einer \EE{psychischen Belastungsreaktion}
bzw. Folge einer Panikattacke. Auch im Rahmen von
(Pop-)Konzerten und ähnlichen Veranstaltungen sind oft
Patienten mit Hyperventilationstetanie anzutreffen.
\bigskip
}{
\begin{itemize}
  \item Tiefe, schnelle Atmung bei sonst gesundem Menschen
  \item Blut-pH wird basisch $\rightarrow$
  Elektrolytstörung
  \item Meist psychisch bedingt (Psych. Belastungsreaktion, Panikattacke, Prüfungen, Veranstaltungen)
%   \begin{itemize}
%     \item 
%     %   \item Lungenembolie
%     %   \item diabetische Stoffwechselentgleisung (Ketoazidose, kompensatorisch)
%   \end{itemize}
\end{itemize}
}{}{}{}

\Fazit{Beim Hyperventilationssysdrom ist die Atmung \EE{ohne 
körperliche Ursache} gesteigert!}



% 1-E
\TextSlide
{\TxABCDE: Hyperventilationssyndrom}
{~
\begin{ABCDEText}
  \item[\ABCDE{1}] Oft kommt es bei vornehmlich
  jugendlichen Besuchern von Pop-Konzerten u.\,ä. zu
  Erregungszuständen, welche zum Hyperventilationssysdrom
  führen können. Daneben kann jedoch jede Situation, welche
  zu erhöhtem Stress oder Erregung führt, ebenso ein
  entsprechender Auslöser sein, 
  siehe \SAMPLER{E}. Das Schaffen bzw. Aufsuchen einer ruhigen
  Umgebung ist wichtig!
  \item[\ABCDE{2}] Häufig kann man Zeichen psychischer
  Erregtheit wahrnehmen: Wildes gestikulieren, Weinkrämpfe
  usw.
  
  Klassisch ist die sog. \T{Pfötchenstellung}:
  Die Arme sind angezogen und Handhaltung erinnert an Pfoten eines Tieres.
  \item[\ABCDE{3}] Wird die Hyperventilation nicht
  rechtzeitig beendet, kann es zur Eintrübung bis hin zur Bewusstlosigkeit kommen.
  \item[\ABCDE{4}] Das Leitsymptom ist die \EE{Atemnot}.
  
  Durch
  das Elektrolyt-Ungleichgewicht kommt es zuerst zu
  \EE{Schwindelgefühlen}, etwas später zu einem
  \EE{Kribbeln} in den
  Fingern. In weiterer Folge zeigen sich tonische \EE{Krämpfe}
  (\T{Hyperventilationstetanie}).
  %   \item[\ABCDE{A}]
  \item[\ABCDE{B}]  Die Patienten beklagen \EE{Atemnot}, obwohl die
  \ce{O2}-Zufuhr nicht beeinträchtigt und das Blut maximal
  gesättigt ist (Sp\ce{O2} von 100\PC*). Sie haben eine
  \EE{tiefe und schnelle Atmung} (\E{Tachypnoe}).
  \item[\ABCDE{C}] Herzfrequenz und Blutdruck können in
  Folge der Erregung etwas erhöht sein.
  %   \item[\ABCDE{D}]
  \item[\ABCDE{E}] \EE{Ausschließen von gefährlichen
  Differentialdiagnosen}: Die psychisch bedingte
  Hyperventilation muß unbedingt von der Hyperventilation 
  aufgrund anderer Ursachen abgegrenzt
  werden. Bei \EE{Ateminsuffizienz} oder \EE{durch
  Erkrankungen verursachte Atemnot} kommt es natürlich auch
  oft zu einer (kompensatorischen) Hyperventilation
  (Herzinfarkt, Lungenembolie, Pneumothorax, Asthma-Anfall,
  \ldots).
  
  Bei einer stoffwechselbedingten \EE{Übersäuerung} (Azidose)
  kommt es zu einer Hyperventilation um \ce{CO2} abzuatmen
  (dadurch wird Säure im Körper abgebaut). Die
  \E{Kussmaul'sche Atmung} beim \EE{diabetischen Koma} wäre
  ein typisches Beispiel dafür
  (\prettyref{topic:koma-diabetisches}).
  \item[\ABCDE{=}] Normalerweise keine vitale Bedrohung
  (Differentialdiagnosen bedenken!).
  %   \item[\ABCDE{\ldots}]
\end{ABCDEText}
}
{
\begin{ABCDESlide}
  \item[\ABCDE{1}] Konzerte, stressige bzw. erregende
  Situationen
  \item[\ABCDE{2}] Gestikulieren, Weinkrämpfe,
  Pfötchenstellung
  \item[\ABCDE{3}] Evtl. Bewusstseinsstörungen
  \item[\ABCDE{4}] Atemnot. Evtl. Kribbeln, Schwindel,
  Krämpfe
  %   \item[\ABCDE{A}]
  \item[\ABCDE{B}] Klagt über Atemnot, Tachypnoe, aber sonst
  keine physischen Anzeichen (keine Zyanose, Sp\ce{O2}
  \SymbolFindingNormal, \ldots)
  \item[\ABCDE{C}] HF und RR evtl. $\uparrow$
  %   \item[\ABCDE{D}]
  \item[\ABCDE{E}] Ausschließen:
  \begin{itemize}
    \item Atemnot mit körperlicher Ursache
    \item Ateminsuffizienz
    \item Stoffwechselbedingte Azidose (z.\,B. Diabetisches Koma)
  \end{itemize}
  \item[\ABCDE{=}] Normal keine vitale Bedrohung
  (Differentialdiagnosen bedenken!).
  %   \item[\ABCDE{\ldots}]
\end{ABCDESlide}
}
{}
{}
{}

\TextSlide
{{\TxSAMPLER}: Hyperventilationssyndrom}
{
\begin{SAMPLERText}
%   \item[\SAMPLER{S}] 
%   \item[\SAMPLER{A}] 
%   \item[\SAMPLER{M}] 
%   \item[\SAMPLER{P}] 
%   \item[\SAMPLER{L}] 
  \item[\SAMPLER{E}] Oft wird über stressige oder aufregende
  Erlebnisse berichtet (Trennung, Todesfälle, Prüfungsstress,
  Verlust des Arbeitsplatzes,
  \ldots).
%   \item[\SAMPLER{R}] 
\end{SAMPLERText}
}
{
\begin{SAMPLERSlide}
%   \item[\SAMPLER{S}] 
%   \item[\SAMPLER{A}] 
%   \item[\SAMPLER{M}] 
%   \item[\SAMPLER{P}] 
%   \item[\SAMPLER{L}] 
  \item[\SAMPLER{E}] Stressige oder aufregende
  Erlebnisse
%   \item[\SAMPLER{R}] 
\end{SAMPLERSlide}
}
{}
{}
{}



% M %%%%%%%%%%%%%%%%%%%%%%%%%%%%%%%%%%%%%%%%%%%%%%%%%%% M %
% Hyperventilationssyndrom
\ProcedurePrintDefault{MF45331C}


\ChapterEnd
